\documentclass[11pt]{article}

\usepackage[margin=1in]{geometry}
\usepackage{amsmath,amssymb,amsthm,mathtools}
\usepackage{microtype}
\usepackage{hyperref}

\hypersetup{
  colorlinks=true,
  linkcolor=blue,
  citecolor=blue,
  urlcolor=blue
}

\newtheorem{theorem}{Theorem}[section]
\newtheorem{lemma}[theorem]{Lemma}
\newtheorem{proposition}[theorem]{Proposition}
\newtheorem{corollary}[theorem]{Corollary}
\theoremstyle{definition}
\newtheorem{definition}[theorem]{Definition}
\newtheorem{remark}[theorem]{Remark}

\newcommand{\R}{\mathbb{R}}
\newcommand{\K}{\mathcal{K}}
\newcommand{\Pcal}{\mathcal{P}}
\newcommand{\E}{\mathbb{E}}
\newcommand{\ip}[2]{\left\langle #1,#2\right\rangle}
\newcommand{\norm}[1]{\left\lVert #1\right\rVert}
\newcommand{\op}{\mathrm{op}}

\title{Structured Spectral Lower Bounds for Quadratic Optimization}
\author{Dmitriy Drusvyatskiy}
\date{\today}

\begin{document}

\maketitle

\begin{abstract}
For convex quadratic minimization, conjugate gradient and more general
Krylov methods are governed exactly by a residual polynomial problem against
the spectral error measure
\[
  \mu_{\mathrm{err}}=\sum_i c_i^2\,\delta_{\lambda_i},
\]
where \(\lambda_i\) are Hessian eigenvalues and \(c_i\) are the coordinates
of the initial error in the eigenbasis.  Classical worst-case lower bounds
show that no first-order method improves the Chebyshev/Krylov rate uniformly
over all smooth quadratics.  In this note we record a sharper structured
version: for any prescribed least-squares spectral error measure, the same
polynomial functional that gives the Krylov upper bound also yields a lower
bound for arbitrary deterministic first-order methods, up to the standard
degree shift \(k\mapsto 2k+1\) from the rotation argument.  The proof combines
three classical ingredients: Gauss quadrature, Jacobi matrices, and the
zero-chain lower-bound construction of Nemirovski and Yudin.
\end{abstract}

\section{Introduction}

Consider the convex quadratic
\[
  f(x)=\frac12 x^\top A x-b^\top x,
  \qquad A=A^\top\succeq 0,
\]
initialized at \(x_0\), and let \(x^\star\) be a minimizer.  When
\(A v_i=\lambda_i v_i\) and
\[
  x_0-x^\star=\sum_i c_i v_i,
\]
the error of conjugate gradient after \(k\) steps is exactly
\[
  f(x_k^{\mathrm{CG}})-f^\star
  =
  \frac12
  \min_{\substack{p\in\Pcal_k\\p(0)=1}}
  \int_0^\beta \lambda p(\lambda)^2\,d\mu_{\mathrm{err}}(\lambda),
  \qquad
  \mu_{\mathrm{err}}:=\sum_i c_i^2\delta_{\lambda_i}.
  \tag{1}
\]
This expression depends not only on the eigenvalue distribution of \(A\), but
also on how the initial error aligns with the eigenbasis.  In least-squares
problems this is the natural measure that appears in the spectral-integral
analysis of gradient descent and conjugate gradient.

The usual worst-case lower bound asks for the slowest possible spectrum and
initial error.  A more refined question is whether a faster rate predicted by
a structured spectral error measure is an artifact of the Krylov restriction,
or whether it is unavoidable for all deterministic first-order methods.  We
show that the structured rate is unavoidable: after a universal degree shift,
the same residual polynomial functional controls the lower bound.

The statement below should be read as a synthesis of standard facts rather
than as a new isolated mechanism.  Gauss quadrature realizes a target measure
by a finite atomic measure with matching polynomial moments; the Jacobi matrix
of this quadrature rule turns the measure into a tridiagonal quadratic hard
instance; and the zero-chain rotation lemma forces any deterministic
first-order method to expose only a growing coordinate flag, which is exactly
the Krylov flag of the Jacobi instance.

\section{Preliminaries}

\subsection{Krylov polynomial identity}

Let
\[
  \K_m(A,r_0)=\operatorname{span}\{r_0,Ar_0,\ldots,A^{m-1}r_0\},
  \qquad r_0=b-Ax_0.
\]
For a quadratic, any point in \(x_0+\K_m(A,r_0)\) can be written as
\[
  x-x^\star=p(A)(x_0-x^\star),
  \qquad p\in\Pcal_m,\quad p(0)=1.
\]
Therefore the best Krylov iterate satisfies the exact identity
\[
  \min_{x\in x_0+\K_m(A,r_0)} f(x)-f^\star
  =
  \frac12
  \min_{\substack{p\in\Pcal_m\\p(0)=1}}
  \int \lambda p(\lambda)^2\,d\mu_{\mathrm{err}}(\lambda).
  \tag{2}
\]

\subsection{Zero-chain quadratics and the rotation lemma}

Let \(E_m=\operatorname{span}\{e_1,\ldots,e_m\}\), with \(E_0=\{0\}\).

\begin{definition}[Zero-chain quadratic]
A convex quadratic \(\bar f:\R^N\to\R\) is a zero-chain quadratic if
\[
  z\in E_m \quad\Longrightarrow\quad \nabla \bar f(z)\in E_{m+1},
  \qquad m=0,\ldots,N-1.
\]
\end{definition}

The following form of the rotation lemma is the standard tool converting
zero-chain structure into lower bounds for arbitrary deterministic first-order
algorithms.

\begin{lemma}[Rotation lemma]
\label{lem:rotation}
Let \(\bar f\) be a zero-chain quadratic on \(\R^N\), fix \(k\ge 0\), and
assume \(N\ge 2k+2\).  For every deterministic first-order algorithm
\(\mathcal A\), there exists an orthogonal matrix \(Q\) such that, when
\(\mathcal A\) is run from \(x_0=0\) on
\[
  f(x)=\bar f(Q^\top x),
\]
the iterates satisfy
\[
  Q^\top x_t\in E_{2t+1},
  \qquad t=0,\ldots,k.
\]
\end{lemma}

The factor \(2t+1\) comes from the adaptive choice of the rotation: at each
oracle query, the construction commits only the directions needed to preserve
past answers while allowing at most two new coordinates to be revealed.

\section{Quadrature and Jacobi Realization}

\subsection{Gauss quadrature}

\begin{lemma}[Gauss quadrature]
\label{lem:gauss}
Let \(\mu\) be a positive Borel measure on \([0,\beta]\), supported on at
least \(N+1\) distinct points, with finite moments up to order \(2N-1\).
Then there exist nodes
\[
  \theta_1<\cdots<\theta_N
\]
in the interior of the convex hull of \(\operatorname{supp}(\mu)\) and
positive weights \(w_1,\ldots,w_N>0\) such that the atomic measure
\[
  \mu_N:=\sum_{j=1}^N w_j\delta_{\theta_j}
\]
satisfies
\[
  \int P\,d\mu_N=\int P\,d\mu
  \qquad
  \forall P\in\Pcal_{2N-1}.
  \tag{3}
\]
\end{lemma}

\begin{proof}
Apply Gram--Schmidt in \(L^2(\mu)\) to the monomial basis
\(1,\lambda,\ldots,\lambda^N\).  This produces orthonormal polynomials
\(\tilde p_0,\ldots,\tilde p_N\) with \(\deg \tilde p_n=n\).  The support
condition ensures that no nonzero polynomial of degree at most \(N\) has zero
\(L^2(\mu)\)-norm.

The degree-\(N\) orthogonal polynomial \(\tilde p_N\) has \(N\) distinct real
zeros \(\theta_1,\ldots,\theta_N\) in the interior of the convex hull of the
support.  Indeed, if it had fewer than \(N\) sign changes on the support, a
polynomial \(r\in\Pcal_{N-1}\) could be chosen so that \(\tilde p_N r\ge 0\)
on the support and is not identically zero there, contradicting
\(\tilde p_N\perp \Pcal_{N-1}\).

Let \(\ell_j\in\Pcal_{N-1}\) be the Lagrange basis at the nodes:
\[
  \ell_j(\theta_i)=\mathbf 1_{i=j}.
\]
Define
\[
  w_j:=\int \ell_j\,d\mu.
\]
For \(P\in\Pcal_{2N-1}\), divide by \(\tilde p_N\):
\[
  P=\tilde p_N s+r,
  \qquad \deg s,\deg r\le N-1.
\]
The quotient term integrates to zero against both \(\mu\), by orthogonality,
and \(\mu_N\), because \(\tilde p_N\) vanishes at every node.  It remains to
integrate \(r\).  Since \(r\in\Pcal_{N-1}\),
\[
  r=\sum_{j=1}^N r(\theta_j)\ell_j,
\]
and therefore
\[
  \int r\,d\mu
  =
  \sum_{j=1}^N r(\theta_j)w_j
  =
  \int r\,d\mu_N.
\]
Finally, \(w_j>0\) because \(\ell_j-\ell_j^2\) vanishes at all quadrature
nodes, hence is divisible by \(\tilde p_N\), so
\[
  w_j=\int \ell_j\,d\mu=\int \ell_j^2\,d\mu>0.
\]
\end{proof}

\subsection{Jacobi realization}

We now turn the finite measure from Gauss quadrature into an actual quadratic
hard instance.  A symmetric tridiagonal matrix is called irreducible if all
off-diagonal entries are nonzero:
\[
  A_{i,i+1}=A_{i+1,i}\ne 0,\qquad i=1,\ldots,N-1.
\]

\begin{lemma}[Jacobi realization]
\label{lem:jacobi}
Let \(\mu\) be a positive Borel measure on \([0,\beta]\) with
\[
  M_2:=\int \lambda^2\,d\mu<\infty,
\]
supported on at least \(N+1:=2k+3\) distinct points in \((0,\beta]\).  Then
there exist an irreducible SPD tridiagonal matrix \(A\in\R^{N\times N}\)
with \(\norm{A}_{\op}\le\beta\) and a scalar \(\tau>0\) such that, for
\[
  b:=\tau e_1,
  \qquad x^\star:=A^{-1}b,
\]
the following holds.  If
\[
  A=\sum_{i=1}^N \lambda_i v_i v_i^\top,
  \qquad
  x^\star=\sum_{i=1}^N c_i v_i,
\]
then the spectral error measure
\[
  \mu_{A,x^\star}:=\sum_{i=1}^N c_i^2\delta_{\lambda_i}
\]
agrees with \(\mu\) on \(\Pcal_{4k+3}\).
\end{lemma}

\begin{proof}
Let \(N=2k+2\), and let \(\mu_N\) be the \(N\)-point Gauss quadrature rule
for \(\mu\) from Lemma~\ref{lem:gauss}.  Since quadrature is exact for
\(\lambda^2\),
\[
  \int \lambda^2\,d\mu_N=M_2.
\]
Define the probability measure
\[
  \nu_N:=\frac{\lambda^2}{M_2}\,\mu_N.
\]

Apply Gram--Schmidt in \(L^2(\nu_N)\) to
\(\{1,\lambda,\ldots,\lambda^{N-1}\}\), obtaining orthonormal polynomials
\[
  q_0\equiv 1,q_1,\ldots,q_{N-1}.
\]
Since \(\nu_N\) has \(N\) atoms, these polynomials form an orthonormal basis
for \(L^2(\nu_N)\).  Let \(A\) be the matrix of multiplication by \(\lambda\)
in this basis:
\[
  A_{ij}:=\ip{q_{i-1}}{\lambda q_{j-1}}_{L^2(\nu_N)}.
\]
Then \(A\) is symmetric and tridiagonal by degree orthogonality.  Its
off-diagonal entries are positive: if \(\kappa_n>0\) is the leading
coefficient of \(q_n\), then comparing leading terms gives
\[
  \lambda q_n
  =
  \frac{\kappa_n}{\kappa_{n+1}}q_{n+1}
  +\text{lower-degree terms},
\]
so
\[
  \ip{q_{n+1}}{\lambda q_n}_{L^2(\nu_N)}
  =
  \frac{\kappa_n}{\kappa_{n+1}}>0.
\]
Thus \(A\) is irreducible.

The spectral bounds follow from the quadratic form.  If
\[
  p_u:=\sum_{n=0}^{N-1}u_{n+1}q_n,
\]
then
\[
  u^\top A u=\int \lambda p_u(\lambda)^2\,d\nu_N(\lambda).
\]
Since \(\nu_N\) is supported on \((0,\beta]\),
\[
  0<u^\top A u
  \le
  \beta \int p_u(\lambda)^2\,d\nu_N(\lambda)
  =
  \beta\norm{u}^2
\]
for every nonzero \(u\).  Hence \(A\) is SPD and \(\norm{A}_{\op}\le\beta\).

Set \(\tau:=\sqrt{M_2}\), \(b:=\tau e_1\), and
\[
  x^\star=A^{-1}b=\tau A^{-1}e_1.
\]
For any polynomial \(P\), the identification of \(A\) with multiplication by
\(\lambda\), together with \(e_1\leftrightarrow q_0\equiv 1\), gives
\[
\begin{aligned}
  x^{\star\top}P(A)x^\star
  &=
  \tau^2 e_1^\top A^{-1}P(A)A^{-1}e_1  \\
  &=
  \tau^2
  \ip{1}{\lambda^{-1}P(\lambda)\lambda^{-1}\cdot 1}_{L^2(\nu_N)}  \\
  &=
  \tau^2\int \lambda^{-2}P(\lambda)\,d\nu_N(\lambda).
\end{aligned}
\]
Thus the spectral error measure of \(x^\star\) is
\[
  \mu_{A,x^\star}
  =
  \tau^2\lambda^{-2}\nu_N
  =
  M_2\lambda^{-2}\left(\frac{\lambda^2}{M_2}\mu_N\right)
  =
  \mu_N.
\]
By Gauss quadrature, \(\mu_N\) agrees with \(\mu\) on
\(\Pcal_{2N-1}=\Pcal_{4k+3}\).
\end{proof}

\section{Main Result}

\begin{theorem}[Structured spectral-error lower bound]
\label{thm:main}
Fix target least-squares spectral data: real numbers
\[
  \lambda_i\in(0,\beta],
  \qquad c_i\in\R,
\]
with at least \(2k+3\) distinct values among those \(\lambda_i\) for which
\(c_i\ne 0\).  Define the target spectral error measure
\[
  \mu_{\mathrm{err}}:=\sum_i c_i^2\delta_{\lambda_i}.
\]
For every deterministic first-order algorithm \(\mathcal A\), there exists a
convex quadratic least-squares instance \(f\) on \(\R^{2k+2}\), initialized at
\(x_0=0\), with \(\norm{\nabla^2 f}_{\op}\le\beta\), whose spectral error
measure agrees with \(\mu_{\mathrm{err}}\) on \(\Pcal_{4k+3}\), such that the
iterate \(x_k\) produced by \(\mathcal A\) satisfies
\[
  f(x_k)-f^\star
  \ge
  \frac12
  \min_{\substack{p\in\Pcal_{2k+1}\\p(0)=1}}
  \int_0^\beta \lambda p(\lambda)^2\,d\mu_{\mathrm{err}}(\lambda).
\]
\end{theorem}

\begin{proof}
Let \(N=2k+2\), and apply Lemma~\ref{lem:jacobi} to
\(\mu_{\mathrm{err}}\), obtaining an irreducible tridiagonal \(A\) and
\(b=\tau e_1\).  Set
\[
  \bar f(z)=\frac12 z^\top A z-b^\top z.
\]
This quadratic is also a least-squares objective: since \(A\succ 0\), choose
\(D=A^{1/2}\) and \(y=A^{-1/2}b\), so that
\[
  \frac12\norm{Dz-y}^2
  =
  \frac12 z^\top A z-b^\top z+\text{constant}.
\]

The quadratic \(\bar f\) is zero-chain.  Indeed, if \(z\in E_m\), then
tridiagonality gives \(Az\in E_{m+1}\), while \(b=\tau e_1\in E_{m+1}\).
Therefore
\[
  \nabla\bar f(z)=Az-b\in E_{m+1}.
\]
By Lemma~\ref{lem:rotation}, there is an orthogonal \(Q\) such that, for
\[
  f(x)=\bar f(Q^\top x),
\]
the iterates of \(\mathcal A\) obey
\[
  Q^\top x_t\in E_{2t+1},\qquad t=0,\ldots,k.
\]

For the Jacobi instance, irreducibility gives
\[
  E_m=\K_m(A,b),\qquad m\le N.
\]
Indeed, tridiagonality gives \(\K_m(A,b)\subseteq E_m\), while the nonzero
subdiagonal entries ensure that \(A^{j-1}e_1\) has a nonzero \(j\)th
coordinate.  Thus \(Q^\top x_k\in\K_{2k+1}(A,b)\), and
\[
\begin{aligned}
  f(x_k)-f^\star
  &=
  \bar f(Q^\top x_k)-\bar f^\star  \\
  &\ge
  \min_{u\in\K_{2k+1}(A,b)} \bar f(u)-\bar f^\star.
\end{aligned}
\]
Using the Krylov polynomial identity \((2)\),
\[
  \min_{u\in\K_{2k+1}(A,b)} \bar f(u)-\bar f^\star
  =
  \frac12
  \min_{\substack{p\in\Pcal_{2k+1}\\p(0)=1}}
  \int_0^\beta \lambda p(\lambda)^2\,d\mu_{A,x^\star}(\lambda).
\]
The integrand has degree at most \(4k+3\), and
\(\mu_{A,x^\star}\) agrees with \(\mu_{\mathrm{err}}\) on
\(\Pcal_{4k+3}\) by Lemma~\ref{lem:jacobi}.  This proves the theorem.
\end{proof}

\begin{corollary}[Krylov rate optimality for prescribed spectral error]
For a prescribed spectral error measure \(\mu_{\mathrm{err}}\), the Krylov
upper bound at degree \(k\) and the arbitrary-first-order lower bound at
degree \(2k+1\) are governed by the same functional
\[
  \mathcal E_m(\mu_{\mathrm{err}})
  :=
  \frac12
  \min_{\substack{p\in\Pcal_m\\p(0)=1}}
  \int \lambda p(\lambda)^2\,d\mu_{\mathrm{err}}(\lambda).
\]
Consequently, any structured rate derived from this functional for Krylov
methods is optimal up to the universal degree shift \(k\leftrightarrow 2k+1\).
\end{corollary}

\section{Examples and Consequences}

The theorem applies directly to the spectral error densities used in refined
least-squares analyses.

\paragraph{Power laws.}
If the spectral error density satisfies
\[
  d\mu_{\mathrm{err}}(\lambda)\approx M\lambda^{a-1}\,d\lambda
  \qquad (a>-1),
\]
then the Krylov polynomial problem yields the familiar rate
\[
  \mathcal E_k(\mu_{\mathrm{err}})
  \asymp k^{-2(a+1)}.
\]
Theorem~\ref{thm:main} implies that no deterministic first-order method can improve
this rate uniformly over least-squares instances with the same spectral error
structure, except for constants and the degree shift \(k\mapsto 2k+1\).

\paragraph{Marchenko--Pastur spectra.}
For random least-squares designs in proportional asymptotics, the reweighted
empirical spectral error measure often converges to a Marchenko--Pastur-type
limit.  The polynomial functional above gives the corresponding Krylov rate
(for example \(k^{-3/2}\) at the hard edge in the square case).  The theorem
shows that this rate is not merely a limitation of Krylov methods; it is a
lower bound for arbitrary deterministic first-order methods on matching
structured instances.

\section{Discussion}

The lower bound depends on the spectral error measure
\(\mu_{\mathrm{err}}\), not on the eigenvalue distribution alone.  Two
least-squares instances may have the same Hessian eigenvalues but different
alignments of the initial error with the eigenvectors, leading to different
polynomial approximation problems and different rates.  This is precisely the
spectral-error-density viewpoint: convergence is controlled by the joint
distribution of eigenvalues and error coordinates.

The exact formulation in Theorem~\ref{thm:main} is best viewed as a derived result
from classical ingredients rather than a standard named theorem in this
generality.  The ingredients are: the zero-chain rotation lower-bound
construction of Nemirovski and Yudin, the residual-polynomial
characterization of Krylov methods and conjugate gradient, and the Jacobi
matrix realization of Gauss quadrature.

\begin{thebibliography}{9}

\bibitem{HestenesStiefel1952}
M.~R. Hestenes and E.~Stiefel.
\newblock Methods of conjugate gradients for solving linear systems.
\newblock \emph{Journal of Research of the National Bureau of Standards},
49(6):409--436, 1952.

\bibitem{Lanczos1952}
C.~Lanczos.
\newblock Solution of systems of linear equations by minimized iterations.
\newblock \emph{Journal of Research of the National Bureau of Standards},
49(1):33--53, 1952.

\bibitem{NemirovskiYudin1983}
A.~S. Nemirovski and D.~B. Yudin.
\newblock \emph{Problem Complexity and Method Efficiency in Optimization}.
\newblock Wiley, 1983.

\bibitem{Nesterov2018}
Y.~Nesterov.
\newblock \emph{Lectures on Convex Optimization}.
\newblock Springer, 2018.

\bibitem{GolubWelsch1969}
G.~H. Golub and J.~H. Welsch.
\newblock Calculation of Gauss quadrature rules.
\newblock \emph{Mathematics of Computation}, 23(106):221--230, 1969.

\bibitem{Gautschi2004}
W.~Gautschi.
\newblock \emph{Orthogonal Polynomials: Computation and Approximation}.
\newblock Oxford University Press, 2004.

\bibitem{Saad2003}
Y.~Saad.
\newblock \emph{Iterative Methods for Sparse Linear Systems}.
\newblock SIAM, 2nd edition, 2003.

\end{thebibliography}

\end{document}
